% Part: normal-modal-logic
% Chapter: filtrations
% Section: S5-decidable

\documentclass[../../../include/open-logic-section]{subfiles}

\begin{document}

\olfileid{nml}{fil}{dec}

\olsection{\Log{S5} সিদ্ধান্তযোগ্য}
প্রমাণগুলির কার্যকর তালিকা এবং সসীম অভীষ্ট ফ্রেমগুলিকে চেনার কার্যকর পরীক্ষা থাকলে, সসীম মডেল ধর্ম দিয়ে স্কিমা-নির্দিষ্ট মোডাল যুক্তির পদ্ধতিগুলি \emph{সিদ্ধান্তযোগ্য} দেখানো যায়; অর্থাৎ, কোনো !!a{formula} ওই পদ্ধতিতে !!{derivable} কি না নির্ণয়ের গণনাযোগ্য পদ্ধতি পাওয়া যায়।
% BN-SRC-722 TARGET-CORRECTION-BEGIN
\emph{উৎস-সংশোধন:} কেবল স্কিমা দিয়ে উল্লেখ করা ও সসীম মডেল ধর্ম থাকা থেকে এই সমান্তরাল গণনা চলে না; প্রমাণ ও অভীষ্ট সসীম প্রতিমডেলের কার্যকর অনুসন্ধানও দরকার। নিচে S5-এর সসীম স্বতঃসিদ্ধ-স্কিমা এবং সর্বজনীন সসীম ফ্রেমে এই শর্তগুলি পূরণ হয়। এই সমান্তরাল প্রমাণে পূর্বনির্ধারিত আকারসীমা লাগে না।
% BN-SRC-722 TARGET-CORRECTION-END


\begin{thm}
  \Log{S5} সিদ্ধান্তযোগ্য।
\end{thm}

\begin{proof}
  % BN-SRC-365: S5-এর প্রতিমডেল খুঁজতে সসীম সর্বজনীন মডেল পরীক্ষা করা হয়।
  ধরা যাক $!A$ দেওয়া আছে এবং $!A$-তে উপস্থিত !!{propositional variable}s $p_1$, \dots, $p_k$-এর মধ্যে পড়ে। প্রতিটি $n$-এর জন্য $n$টি জগৎবিশিষ্ট, $p_1$, \dots, $p_k$-কে মান নির্ধারণকারী মডেলের সংখ্যা সসীম। তাই \emph{সমান্তরালে} দুটি গণনা চালানো যায়: কোনো নিয়মিত পদ্ধতিতে প্রমাণ তৈরি করে \Log{S5}-এর সব উপপাদ্য তালিকাভুক্ত করা; এবং $1$, $2$, \dots জগৎবিশিষ্ট সব সর্বজনীন মডেল তালিকাভুক্ত করে সেগুলির কোনো জগতে $!A$ মিথ্যা কি না পরীক্ষা করা। শেষ পর্যন্ত এই সমান্তরাল প্রক্রিয়া দুটির একটি উত্তর দেবে; কারণ \olref[com][fra]{thm:generaldet} ও \olref[fmp]{cor:S5fmp} অনুসারে হয় $!A$ !!{derivable}, নয়তো কোনো সসীম সর্বজনীন মডেলে এটি মিথ্যা।
\end{proof}

উপরের প্রমাণটি \Log{S5}-এর জন্য কাজ করে, কারণ সর্বজনীন মডেলের ফিল্ট্রেশন আপনা থেকেই সর্বজনীন। প্রতিবিম্বিতা ও সিরিয়ালিটির ক্ষেত্রেও একই কথা সত্য; অন্য ধর্মগুলির জন্য আরও কাজ দরকার।

\end{document}
